\itemheading{6}{คำอธิบายการดำเนินการตาม Model}
วงจรที่เสนอเริ่มจากสำรวจปัญหาและรวบรวมองค์ความรู้ วิเคราะห์กระบวนการและข้อมูล พัฒนาระบบ ทดลองใช้ ติดตามผล ถอดบทเรียน และปรับเป็นแนวปฏิบัติร่วม จากนั้นนำฉบับปรับปรุงกลับไปประเมินซ้ำ โครงสร้างนี้ประยุกต์แนวการเล่าเรื่องวงจรจากตัวอย่าง CPE InternHub ให้เข้ากับงานใบลา โดยไม่ยกผลสำเร็จของโครงการตัวอย่างมาเป็นผลของโครงการนี้

\subsection*{วงรอบที่ 1: ทดลองใช้และค้นหาข้อปรับปรุง}
\textbf{Plan:} จัดทำชุดสถานการณ์ร่วมกับผู้แทนบทบาท ครอบคลุมยื่นใบลา หลายรายวิชา ผู้สอนร่วม ข้อมูลผู้รับผิดชอบไม่ครบ หลักฐานประกอบ สถานะคำขอ และการเข้าถึงข้อมูล พร้อมให้หน่วยงานยืนยันกฎและวิธีนับวัน

\textbf{Do:} ทดลองในสภาพแวดล้อมทดสอบด้วยข้อมูลที่ได้รับอนุญาต บันทึกผลที่คาดหวัง ผลที่เกิดจริง เวลา และข้อเสนอแนะ โดยไม่ใช้ข้อมูลจำลองเพื่ออ้างประสิทธิภาพจริง

\textbf{Check:} เปรียบเทียบคำขอ สถานะ ผู้รับผิดชอบ เอกสาร และยอดรายงานกับข้อมูลต้นทาง ตรวจว่าผู้ใช้เห็นเฉพาะข้อมูลที่มีสิทธิ์ และทบทวนความเข้าใจขั้นตอนของผู้ทดลอง

\textbf{Act:} จัดลำดับปัญหา แก้กฎ การเชื่อมข้อมูล และส่วนติดต่อ แล้วปรับคู่มือกับกรณีทดสอบให้ตรงกัน

\noindent\textbf{หลักฐานที่มีและที่ยังขาด:} การตรวจซอร์สพบว่าหน้ารายงานใช้ข้อมูลจำลอง และมีประเด็นต้องตรวจความครบถ้วนข้อมูลกับกฎการทำงาน แต่ยังไม่มีบันทึกกิจกรรมทดลองใช้หรือผลประเมินวงรอบนี้ จึงไม่ถือว่าปิดวงรอบแล้ว

\subsection*{วงรอบที่ 2: ประเมินฉบับปรับปรุงและยืนยันผล}
\textbf{Plan/Do:} นำฉบับแก้ไขมาทดลองซ้ำด้วยสถานการณ์เดิมและกรณีขอบเขตเพิ่มเติม ตรวจการยื่นและพิจารณาพร้อมกัน การเขียนข้อมูลบริบท การนำเข้ารายเทอม และการเชื่อมรายงานจริง

\textbf{Check:} เปรียบเทียบข้อผิดพลาดและเวลาทำงานระหว่างสองรอบ ตรวจยอดใบลาไม่ซ้ำและผลกระทบรายวิชา ตรวจสิทธิ์ข้ามบุคคลและหน่วยงาน พร้อมให้ผู้รับผิดชอบรับรองผล

\textbf{Act:} สรุปสิ่งที่แก้ได้ ปัญหาคงเหลือ และเงื่อนไขใช้งาน จัดทำคู่มือฉบับรับรองและแผนดูแลระบบก่อนขยายผล

\noindent\textbf{สถานะ:} ยังไม่มีผลประเมินซ้ำหรือหลักฐานผลสัมฤทธิ์ที่ยืนยันความสำเร็จ ต้องแนบวันดำเนินการ ผู้เข้าร่วม รุ่นระบบ รายการปัญหา ผลแก้ไข และผู้รับรองของทั้งสองวงรอบก่อนใช้ข้อความว่าดำเนินการครบแล้ว
